\section{Economic Analysis}

\label{sec:economics}
% ============================================================

\subsection{Fee Competition Model}

Sovereign instances compete for bridge volume through fee setting. We model this as a simultaneous-move game.

\begin{definition}[Fee Competition Game]
\label{def:game}
Let $N$ sovereign instances set fees $\mathbf{f} = (f_1, \ldots, f_N)$ where $f_i \in [0, 100]$ bps. Total bridge volume $V$ is distributed among instances based on fees, liquidity depth $L_i$, and user preferences $\theta_i$:
\begin{equation}
  V_i(\mathbf{f}) = V \cdot \frac{L_i \cdot e^{-\alpha f_i + \theta_i}}{\sum_{j=1}^{N} L_j \cdot e^{-\alpha f_j + \theta_j}}
\end{equation}
where $\alpha > 0$ is the fee sensitivity parameter.
\end{definition}

Each instance maximizes revenue:
\begin{equation}
  \pi_i(\mathbf{f}) = V_i(\mathbf{f}) \cdot \frac{f_i}{10{,}000}
\end{equation}

\begin{theorem}[Nash Equilibrium]
\label{thm:nash}
The fee competition game has a unique symmetric Nash equilibrium at:
\begin{equation}
  f_i^* = \frac{1}{\alpha} \quad \text{for all } i
\end{equation}
when $L_i = L_j$ and $\theta_i = \theta_j$ for all $i, j$ (symmetric instances).
\end{theorem}

\begin{proof}
Taking the first-order condition:
\begin{equation}
  \frac{\partial \pi_i}{\partial f_i} = \frac{V_i}{10{,}000} \left(1 - \alpha f_i \cdot \frac{\sum_{j \neq i} L_j e^{-\alpha f_j + \theta_j}}{\sum_j L_j e^{-\alpha f_j + \theta_j}}\right) = 0
\end{equation}
In the symmetric case, $\frac{\sum_{j \neq i}}{\sum_j} = \frac{N-1}{N}$. Solving:
\begin{equation}
  f_i^* = \frac{N}{\alpha(N-1)}
\end{equation}
As $N \to \infty$, $f_i^* \to 1/\alpha$. For practical values ($N \geq 5$, $\alpha \approx 0.05$), the equilibrium fee is approximately 20~bps, well below the 100~bps hard cap.
\end{proof}

\begin{remark}
The hard cap of 100~bps serves as a credible commitment device. Even a monopolist instance cannot extract more than 1\% per bridge operation. This bounds the worst-case cost for users regardless of market structure.
\end{remark}

\subsection{Asymmetric Equilibria}

When instances differ in liquidity depth $L_i$ or brand preference $\theta_i$:

\begin{proposition}[Liquidity Advantage]
\label{prop:liquidity}
An instance with deeper liquidity ($L_i > L_j$) can sustain a higher fee in equilibrium:
\begin{equation}
  f_i^* > f_j^* \iff L_i > L_j
\end{equation}
Deeper liquidity attracts volume despite higher fees because users face lower slippage.
\end{proposition}

\subsection{xLUX Flywheel Dynamics}

On the Lux C-Chain, all protocol fees from bridge, DEX, lending, perpetuals, and NFT AMM flow to xLUX (LiquidLUX) holders:

\begin{equation}
  \text{APY}_{\text{xLUX}} = \frac{\sum_{p \in \text{protocols}} \Fee_p}{\text{totalStaked}_{\text{LUX}}}
\end{equation}

This creates a positive feedback loop:

\begin{enumerate}[nosep]
  \item More bridged assets $\to$ more TVL on Lux chains.
  \item More TVL $\to$ more DEX/lending/perps activity.
  \item More activity $\to$ more protocol fees.
  \item More fees $\to$ higher xLUX APY.
  \item Higher APY $\to$ more LUX staked as xLUX.
  \item More xLUX $\to$ less circulating LUX $\to$ price support.
  \item Higher LUX price $\to$ deeper bridge liquidity.
  \item Deeper liquidity $\to$ more bridged assets.
\end{enumerate}

\begin{definition}[Flywheel Multiplier]
\label{def:flywheel}
The steady-state flywheel multiplier $\mu$ relates bridge fee revenue to total xLUX yield:
\begin{equation}
  \mu = \frac{\text{APY}_{\text{xLUX}}}{\text{APY}_{\text{bridge-only}}} = 1 + \frac{\sum_{p \neq \text{bridge}} \Fee_p}{\Fee_{\text{bridge}}}
\end{equation}
Empirically, $\mu \approx 4$--$8\times$, meaning bridge fees account for 12--25\% of total xLUX yield.
\end{definition}

\subsection{Optimal Staking Ratio}

The equilibrium staking ratio $s^* = \frac{\text{stakedLUX}}{\text{totalLUX}}$ balances xLUX yield against opportunity cost:

\begin{equation}
  s^* = 1 - \frac{r_{\text{alt}}}{\text{APY}_{\text{xLUX}}(s)}
\end{equation}

where $r_{\text{alt}}$ is the risk-adjusted return available elsewhere. As $s \to 1$, $\text{APY}_{\text{xLUX}}$ increases (denominator decreases), so the equilibrium is stable.

% ============================================================
